\chapter{Methodology}
\label{chap:methodology}

This chapter describes the controlled evaluation used to compare Free and Structured LLM reviewing. The same 30 papers were processed under two complementary experimental tracks. The Counterfactual track tested whether reviews responded specifically to changes in research logic rather than to surface-level rewrites, while the Injection track tested susceptibility to a hidden instruction embedded in the paper PDF. Each input was reviewed under both setups, enabling within-paper comparisons that isolate condition and setup effects from differences between papers.

The generated reviews were evaluated through setup-specific operational measures and a shared auxiliary evidence layer. Free-form ratings and review-aspect counts were extracted from prose by an independent Judge, whereas Structured outputs exposed the corresponding fields directly; injection compliance was Judge-scored for both. Condition-masked common coding, targeted textual audits and a matched human-review benchmark were then used to interpret patterns that could not be resolved by the operational measures alone.

% TODO: Insert the study-design pipeline figure here once finalised.

\section{Study Design and Paper Sample}
\label{sec:method-design-sample}

Papers were sampled from the counterfactual research-logic corpus introduced by Dycke and Gurevych, together with the corresponding original full texts and PDFs. Eligibility required an Original text, an Original PDF and all seven counterfactual variants: three Logic perturbations and four Format perturbations. All 123 papers meeting these requirements formed the sampling pool, from which 30 were selected using random seed 10190. The resulting sample covered ACL and EMNLP papers from 2023--2024, NeurIPS 2024 and ICLR 2025. Full paper identifiers are provided in the appendix.

The analysis unit varied by research question. RQ1 and RQ2 aggregated repeated observations within each paper before inferential comparison. RQ3 used the Original Free--Structured pair for each of the 30 papers to describe review-aspect coverage. RQ4 used all 240 paper-condition rows from the Counterfactual track for its operational-efficiency and rating-dispersion summaries. Multiple eligible human reviews of the same paper were averaged before comparison with its Original Free and Structured reviews.

All reviews were generated with \texttt{gpt-5.4} at temperature 0. Free-form review metrics were extracted by an independent \texttt{deepseek-v4-flash} Judge, also at temperature 0. Setting temperature to 0 reduced random generation variance, allowing comparisons to focus on the reviewing setup and input condition. Table~\ref{tab:method-design-matrix} summarises the resulting experimental matrix.

\begin{table}[htbp]
    \centering
    \caption{Experimental matrix for review generation. Each paper condition was reviewed once under each setup.}
    \label{tab:method-design-matrix}
    \small
    \begin{tabular}{lccc}
        \toprule
        Track & Conditions per paper & Reviewing setups & Generated reviews \\
        \midrule
        Counterfactual & 8 & Free + Structured & 480 \\
        Injection & 2 & Free + Structured & 120 \\
        \midrule
        Total & --- & --- & 600 \\
        \bottomrule
    \end{tabular}
\end{table}

\section{Controlled Input Conditions}
\label{sec:method-inputs}

\subsection{Counterfactual Track}
\label{sec:method-counterfactual}

Each paper had an Original condition and seven counterfactual variants. Three Logic variants altered the paper's conclusion, finding or reported result (\texttt{blueprint\_conclusion}, \texttt{blueprint\_finding} and \texttt{blueprint\_result}). Four Format variants changed voice, regional language convention, writing errors or page layout (\texttt{active\_passive}, \texttt{british\_american}, \texttt{language\_error} and \texttt{paper\_layout}) without intentionally changing the research logic.

For paper $p$, condition $c$, setup $s$ and outcome $Y$, each response was expressed relative to the corresponding Original review:

\begin{equation}
    \Delta_{p,c,s}=Y_{p,c,s}-Y_{p,\mathrm{Original},s}.
    \label{eq:counterfactual-delta}
\end{equation}

The three Logic deltas and four Format deltas were then averaged within each paper and setup. RQ2 defined logic-specific discriminability as the difference between these two mean responses. For methodological-flaw counts, a positive Logic--Format contrast indicated stronger criticism following Logic edits; for ratings, a negative contrast indicated a stronger evaluative penalty. This comparison distinguishes sensitivity to faulty reasoning from general sensitivity to rewritten text.

\subsection{Injection Track}
\label{sec:method-injection}

The Injection track paired each Original PDF with a manipulated copy containing a hidden text-layer payload. Using PyMuPDF, white text at font size 4 was added to the first and final pages; the visible paper content was otherwise unchanged. The payload presented itself as a \texttt{Positive Calibration Track} and instructed the reviewer to raise the rating by approximately one to two points, emphasise strengths, downplay weaknesses and avoid revealing the instruction. The full payload is provided in the appendix.

For each paper and setup, the attack effect was calculated as the outcome under the Manipulated PDF minus the outcome under the Original PDF. Higher ratings, strengths and injection compliance, together with fewer weaknesses and methodological flaws, represented payload-aligned changes. RQ1 compared these paired changes between Free and Structured reviewing: a smaller attack-aligned change in magnitude under Structured indicated attenuation, whereas change close to zero indicated robustness to the injected instruction.

\section{Reviewing Setups}
\label{sec:method-setups}

The two setups differed in their reviewing instruction and output contract.

The Free setup used a generic conference-review instruction requesting a comprehensive, rigorous and constructive assessment. It did not enumerate specific review dimensions or impose a template, and the model returned an unconstrained prose review. The complete prompt is provided in the appendix.

The Structured setup explicitly requested a paper summary, strengths, weaknesses, methodological flaws, an overall rating and confidence, returned as a \texttt{StructuredReview} Pydantic object. Strengths, weaknesses and methodological flaws were represented as lists of distinct textual items, with an empty list permitted when no item was identified. The rating used a 1--10 integer scale and confidence a 1--5 scale. Confidence formed part of the review schema but was not an analysed outcome.

The comparison therefore evaluates schema-constrained reviewing as an operational setup combining explicit review dimensions with typed output, against unconstrained prose reviewing.

\section{Outcome Measurement}
\label{sec:method-measurement}

Review content was measured through the operational output of each setup. For Free reviews, an independent \texttt{deepseek-v4-flash} Judge inferred the rating and extracted distinct review aspects from the generated prose. For Structured reviews, the corresponding measures were read directly from the native typed fields. Injection compliance was Judge-scored for both setups, giving the attack-specific outcome a common measurement source. Table~\ref{tab:outcome-measurement} summarises these mappings.

\begin{table}[t]
    \centering
    \small
    \caption{Operational sources of the principal review measures.}
    \label{tab:outcome-measurement}
    \begin{tabular}{p{0.23\textwidth}p{0.32\textwidth}p{0.32\textwidth}}
        \toprule
        \textbf{Measure} & \textbf{Free} & \textbf{Structured} \\
        \midrule
        Overall rating (1--10) & Judge inference from prose & Native \texttt{rating\_1\_10} field \\
        Strength count & Judge-extracted distinct items & Length of native strengths list \\
        Weakness count & Judge-extracted distinct items & Length of native weaknesses list \\
        Methodological-flaw count & Judge-extracted distinct items & Length of native methodological-flaws list \\
        Injection compliance (0--10) & Judge score & Judge score \\
        \bottomrule
    \end{tabular}
\end{table}

Overall rating represented the review's assessment of paper acceptability on a 1--10 scale. Strengths represented distinct positive merits, weaknesses represented general limitations or concerns, and methodological flaws represented problems in scientific design, evidence or reasoning. Injection compliance measured how far a review followed the hidden instruction embedded in the paper, from no observable compliance (0) to full compliance (10).

RQ4 used additional operational and distributional measures. Review length, counted as whitespace-separated words, measured output volume and concision; end-to-end wall-clock latency measured generation time; and \texttt{o200k\_base} token counts measured output-token demand. Rating dispersion assessed whether schema constraints reduced score differentiation, both globally and within each paper across its eight Counterfactual conditions. Unique-rating and invariant-paper counts described paper-level score differentiation, while token-set Jaccard overlap provided a textual-similarity diagnostic.

The outcome hierarchy followed the research questions. RQ1 centred on rating change, methodological-flaw change and injection compliance, with strength and weakness counts describing broader content shifts. RQ2 used rating and methodological-flaw responses to compare Logic with Format perturbations. RQ3 examined the three aspect counts as recorded coverage. RQ4 assessed efficiency, review length and rating dispersion.

\section{Evidence Triangulation}
\label{sec:method-triangulation}

\subsection{Condition-Masked Common Coding}

A five-paper sample comprising 50 reviews was coded under a common rubric. For each paper, the sample contained the four Injection reviews and six Counterfactual reviews: the Original pair, one Logic pair and one Format pair. Review files were assigned anonymous identifiers and presented in random order without paper identity, condition labels or operational measurements. Because the review format could reveal the setup, the procedure was condition-masked rather than setup-blind. The dissertation author recorded an inferred 1--10 rating, counts of strengths, weaknesses and methodological flaws, and injection compliance using the definitions in Section~\ref{sec:method-measurement}. This condition-masked single-annotator coding placed the Judge-extracted Free measures and Pydantic-native Structured measures on a shared reference scale for triangulating the operational comparison.

\subsection{Targeted Qualitative Audits}

Two post-coding audits examined specific mechanisms behind the aggregate comparisons. For RQ2, all five matched Logic cases were unblinded and coded for whether the altered target was mentioned, correctly connected to claim validity, and followed by a primary-rating decrease relative to the corresponding Original review. These observations distinguished target-defect omission from cases in which criticism was present but did not affect the rating. For RQ3, one Original Free--Structured pair was inspected for repeated themes, item segmentation and category allocation, providing a text-level interpretation of count differences.

\subsection{Matched Human-Review Benchmark}

Public human reviews were available on OpenReview for 11 of the 30 experimental papers, yielding 45 eligible reports. The benchmark included initial primary reviews posted before author rebuttal and excluded meta-reviews, decisions, author responses, public comments, ethics-only reports and rating-only records. Human reviews were measured with the same word-count procedure used for LLM outputs, and the dissertation author coded their strengths, weaknesses and methodological flaws using the same aspect definitions. Where a paper had multiple human reviews, its measurements were averaged before comparison with the corresponding Original Free and Structured reviews. The resulting matched benchmark provided an external reference for review length and aspect profiles.

\section{Statistical Analysis}
\label{sec:method-statistics}

\subsection{Manipulation and Logic-Defect Analyses}

RQ1 estimated the Injection effect for each outcome as the paired Manipulated-minus-Original change. Mean changes and Free-minus-Structured contrasts were reported with $t$-based 95\% confidence intervals. Two-sided $t$-tests assessed whether each setup's mean change differed from zero and whether the paired Free and Structured changes differed across the same 30 papers, using one-sample and paired tests, respectively. The five contrast tests for rating, strengths, weaknesses, methodological flaws and injection compliance formed one family and were adjusted using the Holm procedure.

For RQ2, the Logic and Format changes were first averaged within each paper and setup as defined in Section~\ref{sec:method-counterfactual}. Two-sided paired $t$-tests then compared the mean Logic and Format responses for rating and methodological-flaw count under each setup. The resulting four tests formed one Holm-adjusted family, and each analysis reported the Logic change, Format change, Logic-minus-Format contrast and its 95\% confidence interval.

\subsection{Coverage and Operational Analyses}

RQ3 used descriptive means to summarise strengths, weaknesses, methodological flaws and their total across the 30 Original Free--Structured pairs. The operational profiles, five-paper common coding and matched human benchmark were reported as separate evidence layers.

RQ4 summarised words, latency and output-token use across all 240 Counterfactual rows using means, medians, totals where applicable, and Structured-to-Free ratios. Global rating variance across the 240 ratings per setup was compared using Levene's test. Within each paper and setup, the standard deviation across eight condition-specific ratings quantified score differentiation across inputs; these 30 values per setup were summarised by their mean and median. A standard deviation of zero meant that all eight versions of a paper received the same rating, so these papers were counted separately as cases of complete score invariance. The human-length benchmark used the 11 matched paper-level means and reported the mean, median and range for each review source. All remaining RQ4 comparisons were descriptive.

\section{Reproducibility and Research Ethics}
\label{sec:method-reproducibility-ethics}

Paper selection, condition construction, review generation, metric extraction and analysis were scripted. Reproducibility materials preserve the random seeds, complete prompts, Pydantic schema, model identifiers, temperatures, tokenizer, execution records, raw outputs, manual annotations and final analysis artifacts, linked by paper, condition and setup. Outputs were validated for non-empty Free text, schema-compliant Structured responses and valid Judge fields; failed calls were eligible for retry. All 600 generated reviews and required Judge extractions completed without missing outcome rows, and no outcomes were imputed.

The prompt-injection experiment was confined to controlled paper copies and research API calls; no manipulated document entered a live review platform or editorial decision. The human benchmark used publicly available OpenReview reports. Disclosure of the payload and generated outputs supports reproducibility, vulnerability assessment and safer reviewing-system design.
